% !TeX root = main.tex
\section{项目简介}

\subsection{项目定位与核心需求}

OpenWrite 是一套面向网络小说作者、独立长篇创作者及小型编辑团队的 Agent 创作系统。作者提出题材、人物动机和故事承诺后，智能体团队协助整理人物与世界观、构建分层大纲、生成章节、核查连续性并提出修订；作者保留设定定稿、重要改稿与最终交付的决定权。

长篇创作的困难会随章节累积：作者创作难以统筹全书设定，人物认知、伤势、世界观与伏笔容易相互冲突；写作者与审稿者如果共享同一思路，可能遗漏结构性问题；多轮修改后，正文、章节摘要和人物状态又可能出现版本脱节。创作者因此需要反复整理资料和核对改动，编辑也难以判断一份审稿意见究竟针对哪版正文。

本项目将这些需求转化为三个可操作的目标：\textbf{写前有据、审稿有证、改后可核验。}每章写作在统一上下文的前提下出发，审稿意见关联正文证据，修订通过差异预览、版本重验使 Agent 能够承担持续性的创作事务。

\subsection{已有基础与迁移方向}

本作品基于 GitHub 开源项目 \href{https://github.com/LiPu-jpg/Openwrite}{LiPu-jpg/Openwrite} 继续建设。其 DeepSeek Harness 插件版本已经形成统一创作 Agent、小说领域工具桥接、创作工作台和章节审稿查询图等实现；OpenWrite Core 提供独立的小说领域能力，保存人物、设定、大纲、正文及修订记录，并提供上下文装配、评审、接纳与导出能力。本稿撰写之日仓库页面显示 826 Stars 与 127 Forks，体现了已有社区关注。

本次项目拟迁移至 openJiuwen：由 JiuwenSwarm 组织 Agent 团队，以 SwarmFlow 描述阶段顺序、并行评审和修订分支，以 Skills 沉淀角色方法，以 MCP 调用 OpenWrite 既有领域服务。Planner-Agent、Writer-Agent、三名专职 Review-Agent 和 Revision-Agent 构成六个 AI 工作角色，由 Leader-Agent 统筹；作者作为人工确认方参与关键决策。现有迁移分支已包含工作流、角色与 MCP 桥接雏形，本方案将其作为实施基础，继续完成协议对齐和真实运行验收。

\subsection{创新性与实用价值}

\begin{enumerate}
  \item \textbf{面向作品状态的长期协作。}将作者意图、分层大纲、人物时态、章节记忆和伏笔组织成写作依据，Agent 每次执行任务均读取相关信息，降低长对话中遗漏关键事实的风险。
  \item \textbf{写作与评审分工，程序统一判定。} 设计由三个 Review-Agent 独立覆盖六个质量域，并沿用领域的证据；由程序聚合质量分和覆盖率，向作者呈现问题与建议。分工有助于提供不同职责视角，实际检出效果仍需人工评估。
  \item \textbf{将改稿纳入版本闭环。}正文的每次修改关联到具体提案和来源版本；旧评审在改稿后失效，影响到的人物状态和章节记忆须重建或进行复核，避免“文字已改、故事状态未改”的情况发生。
  \item \textbf{沉淀可复用的开源领域能力。}在既有项目基础上，计划开放团队技能、MCP 适配与可复现样例，使社区能够复用创作流程、扩展题材规则，并比较迁移前后的行为；随包第三方内容保留原许可与来源。
\end{enumerate}

项目以可持续的单章闭环作为首要交付单位，逐步扩展到多章连载与旧稿维护。预期价值是减少资料整理和重复核查的工作量，让作者将更多精力用于叙事判断与故事创作；具体效率与质量收益将在真实章节及人工评估中测量。
